\section{A merge tree for a low-density spectrum}
\label{sec:physical_ld_tree}

The interval families of Lemma~\ref{lem:dyadic_interval_low_energy_family}
compare two different low-energy windows: viability for the upper window
of width $\eta$ is carried to viability for the lower window of width
$\eta-3w$. A merge tree built from them therefore uses one window per
level, shortened by the buffer $3w$ at each of its levels, as in the
condition $\eta_1\ge2\mu/\log n$ and the buffer $w=\mu/(3\log n)$
of~\cite[Theorem~\texttt{t:aisgapless}, lines 1027--1041]{Arad2016RigorousRG}.
The statement below is the viability induction of the merge tree
of~\cite[Theorem~\texttt{t:alg2}, lines 1869--1878]{Arad2016RigorousRG}.
It is restricted to long nodes that keep the whole family image of the
merged child spaces: no random sampling, collective trimming, bound on
the node dimensions, final refinement or matrix product state bound of
the output is asserted, and no running time. The constant depends on the
fixed bandwidth, as recorded in~\cite{gap:alvv16_low_energy_family_scope}.

\begin{lemma}[Midpoint induction for the interval recursion]
    \label{lem:balanced_interval_invariant_midpoint}
    \lean{TNLean.MPS.AreaLaw.balancedIntervalSpace_invariant_of_midpoint}
    \leanok
    \uses{def:balanced_interval_spaces}
    Any property of interval spaces that holds at the short leaves, and is
    preserved by the parent operation applied at the actual midpoint
    $k=p+\lfloor(r-p)/2\rfloor$ of a long interval $[p,r)$, holds at every
    recursively constructed space.
\end{lemma}
\begin{proof}\leanok
    Induct on the interval length, as in
    Lemma~\ref{lem:balanced_interval_invariant}, retaining that the cut of
    a long interval is its midpoint.
\end{proof}

\begin{lemma}[Viability at every level of the low-density merge tree]
    \label{lem:physical_ld_interval_tree}
    \lean{TNLean.MPS.AreaLaw.exists_physical_ld_interval_tree}
    \leanok
    \uses{lem:dyadic_interval_low_energy_family,
        lem:balanced_interval_invariant_midpoint,
        thm:actual_adjacent_interval_merge,
        lem:full_interval_viability}
    Fix $\eta_{\rm LD}>0$. There is a constant $C\ge\max\{100,19+\eta_{\rm LD}\}$
    with the following property. Let $q\ge2$, $m\ge q$ and $w>0$, put
    $g=\min\{w,1\}$, $u=1+\log m+\log(1/g)$,
    $\ell=\lceil Cu^2/g\rceil$ and $D=(m\ell^2)^{22\ell}$. Let
    $H=\sum_ih_i$ be a nearest-neighbour chain of length $L$ and local
    dimension $q$ with $0\le h_i\le\Id$, and let $G(\eta)$ be the span of
    the eigenvectors of $H$ with energy at most $E_0+\eta$, where $E_0$ is
    the ground energy. Then $\ell\ge21$, $\log D$ obeys the bound of
    Lemma~\ref{lem:dyadic_interval_low_energy_family}, and there are
    operators $A_{J,1},\dots,A_{J,D^2}$ on every interval $J$ with the
    following property. In the balanced recursion with leaf width
    $4\ell$, let a long node $J=[p,r)$ with children $S_1,S_2$ retain
    \[
        \sum_{i=1}^{D^2}A_{J,i}\,(S_1\otimes S_2).
    \]
    If $r-p\le4\ell\,2^k$ and $3w(k+1)\le\eta_{\rm LD}$, the space on $J$ is
    $3/200$-viable for $G(\eta_{\rm LD}-3wk)$. In particular, if
    $3w(\lceil\log_2L\rceil+1)\le\eta_{\rm LD}$, the space at the root is
    $3/200$-close to $G(\eta_{\rm LD}-3w\lceil\log_2L\rceil)$.
\end{lemma}
\begin{proof}\leanok
    \uses{lem:dyadic_interval_low_energy_family,
        lem:balanced_interval_invariant_midpoint,
        thm:actual_adjacent_interval_merge,
        lem:full_interval_viability}
    Take the family of Lemma~\ref{lem:dyadic_interval_low_energy_family}
    on every interval, at the index zero, so that its comparison loss is
    $1/200$ and its suppression satisfies $\Delta\le D^{12}\Delta\le10^{-5}$.
    Induct over the recursion with
    Lemma~\ref{lem:balanced_interval_invariant_midpoint}. A leaf retains
    the full local space, which is viable for every target. A long node of
    length at most $4\ell\,2^{k+1}$ has children of length at most
    $4\ell\,2^k$, which are $3/200$-viable for $G(\eta)$ with
    $\eta=\eta_{\rm LD}-3wk\ge6w$. By
    Lemma~\ref{thm:actual_adjacent_interval_merge}, their merge is
    $3/100$-viable for $G(\eta)$, hence $7/100$-viable for the upper window
    of the family, which is $1/200$-close to it. The family image is then
    $\Delta/(1-7/100)^2$-viable for the lower window, which is
    $1/200$-close to $G(\eta-3w)$. The resulting error
    $2\bigl(\Delta/(93/100)^2+1/200\bigr)$ is at most $3/200$. At the root,
    $L\le2^{\lceil\log_2L\rceil}$, and
    Lemma~\ref{lem:full_interval_viability} turns viability on the whole
    chain into closeness.
\end{proof}
